%% kapa: exam=1 round=35 year=2024 period=1 form=A subject=civil_law subject_ko=민법 seq=13 number=13
%% source: Q-Net 감정평가사 기출문제 (cst00309) artlSeq=5213229 file=raw/1008/5213229/1교시 문제지 원본.pdf
%% conversion: kordoc 4.24.51 -> tools/gen1.py (auto); review=proofread-sonnet; flags=none
\begin{question}{13}
표현대리에 관한 설명으로 옳지 않은 것은? (다툼이 있으면 판례에 따름)
\begin{choices}
  \item 표현대리행위가 성립하는 경우, 상대방에게 과실이 있더라도 과실상계의 법리를 유추 적용하여 본인의 책임을 경감할 수 없다.
  \item 상대방의 유권대리 주장에는 표현대리의 주장이 포함되는 것은 아니므로 이 경우 법원은 표현대리의 성립여부까지 판단해야 하는 것은 아니다.
  \item 민법 제126조의 권한을 넘은 표현대리 규정은 법정대리에도 적용된다.
  \item 복대리인의 대리행위에 대해서는 표현대리가 성립할 수 없다.
  \item 수권행위가 무효인 경우, 민법 제129조의 대리권 소멸 후의 표현대리가 적용되지 않는다.
\end{choices}
\end{question}
\answer{4}
